\begin{block}{Input data}

\begin{figure}[H]\centering
\includesvg[inkscapelatex=false, width=0.9\textwidth]{../images/fig3-resource-relationship}
\setcounter{figure}{4}
\caption{Relationships among the project's data sources and tools: TEI/XML merge,
vocabulary datasets, search and network-visualization tools, translation corpus.}
\end{figure}

\begin{enumerate}
\item \textbf{Base TEI/XML}: the Kokinwakashu text, five-ku segmentation, author
  information, textual variants, commentary, and other existing TEI markup
\item \textbf{Vocabulary and morphological information}: the Hachidaishu vocabulary
  list and part-of-speech data; used to infer kugire candidates from inflection
  forms, sentence-final particles, binding particles, etc.
\item \textbf{Modern Japanese translation}: the modern Japanese gloss in Kaneko
  Motoomi's \textit{Kokinwakashu Hy\=osh\=aku}; used to observe how 20th-century
  commentators construed the waka as modern Japanese discourse
\end{enumerate}
\end{block}
